\section{Contribuições à Ciência (Científicas, Tecnológicas ou de Inovação)}
% TODOS AUTORES (MÁX. 15) DEVERÃO CONSTAR DA CITAÇÃO
% USAR PADRÃO BIBLIOGRÁFICO DE SUA ÁREA
% JUSTIFICATIVA:
% 1. IMPACTO
% 2. RELEVÂNCIA PARA ESTA PRODUÇÃO
% 3. IMPORTÂNCIA PARA SEUS OUTROS PROJETOS
\begin{enumerate}[label=$\bullet$, itemsep=1.25\baselineskip, topsep=0.5\baselineskip]
	\item
	\fullcite{neto-setti2025}

	\medskip
	\textbf{Justificativa:}
	% 1. IMPACTO
	Esta pesquisa possui impacto científico pois demonstra rigorosamente tanto a convergência do método de Kaczmarz, quanto sua ordem linear.
	% 2. RELEVÂNCIA PARA ESTA PRODUÇÃO
	É de grande relevância a esta produção, pois o método consta dentre os assuntos preliminares para os vídeos prometidos no plano de atividades.
%	% 3. IMPORTÂNCIA PARA SEUS OUTROS PROJETOS
	Por fim, expôs-me à reconstrução tomográfica --- um clássico problema inverso --- como também conceitos de otimização convexa, aprimorando minha capacidade de realizar projetos futuros sobre tais temas.


	\item
	\fullcite{valem2026}

	\medskip
	\textbf{Justificativa:}
	% 1. IMPACTO
	Este projeto possui impacto social pois mitiga a barreira linguística no estudo de IA.
	% 2. RELEVÂNCIA PARA ESTA PRODUÇÃO
	Nele, lecionei sobre \emph{backpropagation} --- algoritmo intimamente ligado a otimização e gradientes, outro tema dentre os do plano de atividades --- e, de forma mais geral, exigiu de mim produzir vídeo-aulas pedagógicas. Evidentemente, tal exigência é de extrema pertinência a esta produção, também centrada em vídeo-aulas,
	% 3. IMPORTÂNCIA PARA SEUS OUTROS PROJETOS
	assim como é basilar para minhas perspectivas futuras de produções pedagógicas.
\end{enumerate}
